\chaptertitle{第六章\quad 统计与概率的综合应用}

前五章我们分别学了数据的收集、整理、描述和分析，也学了概率的基础知识。这一章把它们综合起来，看看这些工具在实际生活中怎么用、怎么避免被数据欺骗——以及统计思想怎么催生了现代人工智能中最基础的预测方法。

\sectiontitle{6.1\quad 用数据做决策}

当面对不确定的情况时，统计和概率帮助我们做出更合理的决策。

\begin{example}
两个水果店，甲店苹果合格率$95\%$，乙店$90\%$，价格相同。在哪家买？
\end{example}
\begin{solution}
合格率甲店更高，价格相同——当然选甲店。但如果甲店价格是乙店的$2$倍呢？那就需要权衡：乙店虽然坏果率高一倍，但价格便宜一半，可能更划算。统计为决策提供了依据，但最终选择还取决于你的价值判断。
\end{solution}

\begin{example}
天气预报降雨概率$80\%$，你会带伞吗？
\end{example}
\begin{solution}
$80\%$意味着类似条件下$80\%$的日子会下雨，概率很高，带伞明智。如果只有$10\%$，很多人就不带了。每个人的风险承受能力不同，但数据提供了理性决策的基础。
\end{solution}

\begin{example}
某公司招聘，笔试通过率$30\%$。小明考后排前$40\%$，他能通过吗？
\end{example}
\begin{solution}
不一定。通过率只有$30\%$，小明排$40\%$——说明他还不够靠前。只有前$30\%$的人能通过，小明可能在淘汰边缘。当然，如果前面有人放弃，他可能递补。数据告诉我们"可能性有多大"，但不能保证结果。
\end{solution}

\begin{exercise}
\item 某药治愈率$85\%$，另一种$70\%$但价格便宜一半。你选哪种？为什么？
\item 甲股票平均年收益$12\%$，标准差$15\%$；乙股票平均$8\%$，标准差$5\%$。风险厌恶者选哪支？
\item 老师习惯叫学号末位$0$或$5$的同学。小明学号$8$，被点到的概率是多少？这个例子说明什么？
\end{exercise}

\sectiontitle{6.2\quad 统计陷阱——如何不被数据欺骗}

"数据不会说谎，但说谎者会用数据。"以下是一些常见的统计陷阱。

\textbf{陷阱一：样本偏差。}

\begin{example}
某调查发现，使用某品牌牙膏的人中$90\%$牙齿更白了。广告："$90\%$的用户牙齿变白！"可信吗？
\end{example}
\begin{solution}
需要追问：样本怎么选的？如果只调查了$10$个人，或者只调查了该品牌的忠实用户，结果就不可靠。样本必须能代表总体。
\end{solution}

\textbf{陷阱二：平均数掩盖差异。}

\begin{example}
某公司公布"员工平均年薪$30$万"，但$90\%$员工年薪不到$10$万，高管人均$200$万。
\end{example}
\begin{solution}
平均数被极端值拉高。这时应该看中位数——如果中位数只有$8$万，就知道"平均$30$万"是误导。
\end{solution}

\textbf{陷阱三：忽略基数。}

\begin{example}
甲医院$100$台手术成功$90$台（成功率$90\%$），乙医院$1000$台成功$890$台（$89\%$）。甲广告说"成功率更高"。
\end{example}
\begin{solution}
乙虽然成功率略低，但做了$10$倍的手术，经验可能更丰富。而且$90\%$和$89\%$的差异可能只是随机波动。只看百分比忽略基数，会得出片面结论。
\end{solution}

\textbf{陷阱四：相关不等于因果。}

\begin{example}
冰淇淋销量越高，溺水事故越多。所以"吃冰淇淋导致溺水"？
\end{example}
\begin{solution}
当然不是！共同原因是"夏天到了"——天热冰淇淋卖得好，游泳的人也多，溺水就多了。两个变量相关（一起变化），不等于有因果关系。这是统计误导中最常见、也最隐蔽的一种。
\end{solution}

\knowbox{
\textbf{四大统计陷阱}：\\
1. \textbf{样本偏差}——样本不代表总体。\\
2. \textbf{平均数陷阱}——极端值拉偏。\\
3. \textbf{忽略基数}——只看比例不看总量。\\
4. \textbf{相关≠因果}——同时发生不等于一个导致另一个。}

\begin{exercise}
\item 指出以下"统计"中的陷阱：\\
(a) "本产品$80\%$用户满意！"（调查了$5$人）\\
(b) "重点高中$95\%$考上大学"（该校本就招收最优秀的学生）\\
(c) "交通事故比前年增$50\%$"（前年$2$起，去年$3$起）\\
(d) "喝咖啡的人心脏病发病率更高"（喝咖啡的人可能更爱熬夜）
\item 某补习班广告："培训后平均成绩提高$20$分！"哪里可能误导？
\item "我穿红色衣服考试得了满分，所以红色衣服带来好运。"怎么用统计学反驳？
\end{exercise}

\sectiontitle{6.3\quad 从统计到机器学习——最简单的二分预测}

你可能听说过"人工智能"、"机器学习"这些词。其实机器学习中最基础的方法——**决策树**——背后的核心思想就是统计：用已有数据发现规律，然后用这些规律做预测。

来看一个最简单的例子。假设你收集了一些关于"周末是否去公园"的数据：

\begin{center}
\begin{tabular}{c|c|c|c}
天气 & 温度 & 是否节假日 & 去公园？\\\hline
晴 & 高 & 是 & 去\\
晴 & 高 & 否 & 去\\
阴 & 高 & 是 & 去\\
雨 & 高 & 是 & 不去\\
雨 & 低 & 是 & 不去\\
雨 & 低 & 否 & 不去\\
阴 & 低 & 是 & 去\\
晴 & 低 & 否 & 去
\end{tabular}
\end{center}

现在有个问题：如果\textbf{天气=雨、温度=高、是否节假日=否}，你会去公园吗？

凭直觉，你可能就猜"不去"——因为已有的数据中，下雨时从来不去。但我们需要一个系统的方法。统计学家会怎么做？看已有的数据中，每种情况出现了多少次，然后算概率。

\begin{example}
根据上表数据，计算以下概率：\\
(1) 总体中去公园的概率；(2) 下雨天去公园的概率。
\end{example}
\begin{solution}
(1) 总共$8$条数据，去公园$5$次，$P(\text{去})=\dfrac58=0.625$。\\
(2) 下雨的天数有$3$条（(雨,高,是)、(雨,低,是)、(雨,低,否)），其中去公园$0$次。$P(\text{去}|\text{下雨})=\dfrac03=0$。\\
所以如果已知下雨，去公园的概率为$0$——这就是根据数据做出的预测！
\end{solution}

这个思路就是\textbf{决策树}的雏形：先找一个"最佳特征"来分割数据（比如先看天气：如果下雨就直接预测"不去"；如果晴或阴，再根据其他特征判断），不断分割下去，直到每个分支的数据足够"纯"（全是去或全是不去）。

\begin{center}
\begin{tikzpicture}[scale=0.8]
\node[draw,thick,fill=lightblue] (root) at (3,4){天气？};
\node[draw,thick,fill=green!20] (sun) at (0,2){晴/阴};
\node[draw,thick,fill=exred!20] (rain) at (6,2){雨};
\node[draw,thick,fill=green!20] (temp) at (0,0){温度？};
\node[draw,thick] (sunny_high) at (-1.5,-2){去};
\node[draw,thick] (sunny_low) at (1.5,-2){去};
\node[draw,thick,fill=exred!20] (no) at (6,0){不去};
\draw[thick] (root)--(sun) node[midway,left]{是};
\draw[thick] (root)--(rain) node[midway,right]{否};
\draw[thick] (sun)--(temp);
\draw[thick] (temp)--(sunny_high) node[midway,left]{高};
\draw[thick] (temp)--(sunny_low) node[midway,right]{低};
\draw[thick] (rain)--(no);
\node at (3,-3.2){\small 图6-1\quad 最简单的决策树："去不去公园？"};
\end{tikzpicture}
\end{center}

当然，现实中很少遇到这么"完美"的数据——所有的下雨天都不去。更多的是：下雨天$80\%$不去、$20\%$去。这时我们就需要进行\textbf{概率预测}："下雨天不去的概率是$80\%$"——这已经足够帮助我们做决策了。

\begin{example}
一家银行想根据客户的资料预测"是否会按时还款"。他们收集了历史数据，发现：
\begin{itemize}
\item 有稳定收入的客户中，$95\%$按时还款；
\item 没有稳定收入但有抵押的客户中，$80\%$按时还款；
\item 既无稳定收入又无抵押的客户中，只有$40\%$按时还款。
\end{itemize}
如果有一个新客户：有稳定收入、无抵押，银行应该批准贷款吗？
\end{example}
\begin{solution}
根据历史数据，有稳定收入的客户按时还款概率$95\%$——很高。银行可以批准。\\
当然$95\%$不是说$100\%$安全，而是说在大量类似客户中，只有$5\%$会违约。银行可以接受这个风险水平。这就是统计在金融决策中的应用。
\end{solution}

\knowbox{
\textbf{机器学习与统计的关系}：\\
机器学习中的许多基础方法（决策树、朴素贝叶斯、线性回归）\\
本质上是统计学的延伸——用历史数据中的频率估计概率，\\
然后用概率做预测。学好了统计与概率，\\
你就已经迈出了理解人工智能的第一步。}

\begin{exercise}
\item 根据"去公园"的数据表，如果天气=晴、温度=低、节假日=是，你会预测去不去？为什么？
\item 某邮箱服务商收集了$1000$封邮件，发现：含"免费"字样的邮件中有$90\%$是垃圾邮件，不含的只有$5\%$是垃圾邮件。现在收到一封含"免费"字样的新邮件，它是垃圾邮件的概率是多少？
\item 除了决策树，你还能想到哪些日常生活中的"根据历史数据做预测"的例子？（提示：电商推荐、天气预报、电影评分）
\end{exercise}

\sectiontitle{6.4\quad 项目：设计一个统计调查}

这一节是一个动手项目，把全书的知识串联起来。

\begin{example}
小明想调查"本校同学最喜欢哪种休闲活动"，设计完整的统计调查方案。
\end{example}
\begin{solution}
完整流程：\\
\textbf{1. 明确问题。}本校同学最喜欢的休闲活动是什么？\\
\textbf{2. 确定调查方式。}全校$2000$人，普查太麻烦，用抽样调查。\\
\textbf{3. 设计抽样方案。}按年级分层，各年级随机抽$50$人，共$150$人。\\
\textbf{4. 设计问卷。}列出常见选项（阅读、运动、游戏、看视频、其他），每人选一项。\\
\textbf{5. 收集数据。}分发问卷并回收。\\
\textbf{6. 整理数据。}制作频数-频率表。\\
\textbf{7. 画统计图。}条形图比大小，扇形图看比例。\\
\textbf{8. 分析数据。}哪项最受欢迎？各年级有差异吗？\\
\textbf{9. 得出结论。}写成调查报告。
\end{solution}

\knowbox{
\textbf{统计调查完整流程}：\\
明确问题$\rightarrow$确定方式$\rightarrow$设计方案$\rightarrow$设计问卷\\
$\rightarrow$收集数据$\rightarrow$整理数据$\rightarrow$画统计图\\
$\rightarrow$分析数据$\rightarrow$得出结论}

\begin{exercise}
\item 设计调查：了解你们班同学平均每天课外阅读时间。按9步写出方案。
\item 全校$2000$人，按年级分层抽样和简单随机抽样各有什么优缺点？
\item 调查报告除了数据和图表，还应包括什么？（至少3点）
\end{exercise}

\sectiontitle{本章小结}

这一章我们把全书内容用到实际中。用数据做决策的核心是：收集可靠数据，用合适的方法分析，然后理性判断。同时要警惕数据陷阱——样本偏差、平均数误导、忽略基数、相关误为因果。我们还看到，机器学习中最基础的决策树方法，本质上就是统计中的条件概率——用历史数据的频率来做预测。最后，一个完整的统计调查项目，是检验你综合运用统计能力的最好方式。

\sectiontitle{课后练习}

\subsection*{A组\quad 基础巩固}

\begin{exercise}
\item 某种病发病率$0.1\%$，检测准确率$99\%$（患者检出阳性$99\%$，健康人检出阴性$99\%$）。一个人检测阳性，他患病的概率是多少？（提示：假设$10000$人，算真阳性和假阳性）
\item 餐厅$50$份问卷中$45$人满意，老板打横幅"$90\%$顾客满意！"有什么问题？
\item 某地区GDP增长和人均寿命正相关。能说"GDP增长导致寿命增长"吗？为什么？
\item 统计图表纵轴不从$0$开始会夸大差异——请画一个简单例子说明。
\item 根据"去公园"数据表：天气=晴、温度=高、非节假日，预测去不去公园，并说明依据。
\end{exercise}

\subsection*{B组\quad 挑战提升}

\begin{exercise}
\item 某学校$600$男生、$400$女生。调查校服满意度：\\
方案A：从$1000$人中随机抽$100$人；方案B：男生抽$60$人、女生抽$40$人。哪个更好？各是什么抽样方法？
\item 电商好评率$98\%$，但只有$5\%$买家写了评价。可信吗？如何改进评价机制？
\item 小明和小红游戏：同时抛两枚硬币，都正面小明得$3$分，都反面小明得$1$分，一正一反小红得$2$分。公平吗？用概率和期望值计算。
\item 调查你所在班级/年级，完成一份关于"最关心的校园问题"的完整调查报告。
\item （选做）找一份真实的新闻报道，其中用到了统计数据。分析它有没有犯本章6.2节中提到的任何一种统计陷阱。
\end{exercise}

\newpage
